\chapter{Prompt Reference}\label{ch:prompt-reference}\index{Prompt Reference}

This appendix provides essential prompts organized by situation. These are not scripts to recite---they're starting points for genuine inquiry. Adapt them to your context.

\section{Clarification Prompts}

Use when requirements or direction are unclear.

\begin{table}[H]
\centering
\begin{tabular}{p{3cm}p{9cm}}
\toprule
\textbf{Situation} & \textbf{Prompt} \\
\midrule
Unclear problem & ``What problem are we actually solving? What does success look like?'' \\
Ambiguous scope & ``What's in scope? What should I explicitly NOT build?'' \\
Multiple interpretations & ``This could mean A or B. Which direction should we go?'' \\
Missing context & ``What prerequisite knowledge am I missing to understand this?'' \\
\bottomrule
\end{tabular}
\end{table}

\section{Challenge Prompts}

Use to surface assumptions and stress-test decisions.

\begin{table}[H]
\centering
\begin{tabular}{p{3cm}p{9cm}}
\toprule
\textbf{Situation} & \textbf{Prompt} \\
\midrule
Question assumptions & ``What assumptions are built into that approach?'' \\
Test robustness & ``What's the strongest argument against this?'' \\
Anticipate failure & ``If this fails, what's the most likely reason?'' \\
Find blind spots & ``What am I not seeing about this problem?'' \\
Future regret & ``What would make us regret this decision in 6 months?'' \\
\bottomrule
\end{tabular}
\end{table}

\section{Plan Review Prompts}

Use at checkpoints before approving AI-generated plans.

\begin{table}[H]
\centering
\begin{tabular}{p{3cm}p{9cm}}
\toprule
\textbf{Situation} & \textbf{Prompt} \\
\midrule
Standard checkpoint & ``Before I approve: What alternatives did you consider? What's the riskiest assumption? What would make this fail?'' \\
Complexity check & ``Is this complexity necessary, or are we over-engineering?'' \\
Trade-off analysis & ``What does this optimize for? What does it sacrifice?'' \\
Simplification & ``What's the simplest version that could work?'' \\
\bottomrule
\end{tabular}
\end{table}

\section{Learning Prompts}

Use when the goal is understanding, not just output.

\begin{table}[H]
\centering
\begin{tabular}{p{3cm}p{9cm}}
\toprule
\textbf{Situation} & \textbf{Prompt} \\
\midrule
Guided discovery & ``Don't give me the answer---ask me questions that help me discover it.'' \\
Key insight & ``What's the key insight that would unlock my understanding?'' \\
Test understanding & ``I think I understand X. Quiz me to find my gaps.'' \\
Build depth & ``What's the next level of depth I should explore?'' \\
\bottomrule
\end{tabular}
\end{table}

\section{Debugging Prompts}

Use when diagnosing problems.

\begin{table}[H]
\centering
\begin{tabular}{p{3cm}p{9cm}}
\toprule
\textbf{Situation} & \textbf{Prompt} \\
\midrule
Systematic diagnosis & ``What are the possible causes? Let's rank them by likelihood.'' \\
Assumption check & ``What assumptions am I making about how this should work?'' \\
Minimal reproduction & ``What's the smallest test case that reproduces this?'' \\
Getting unstuck & ``I've been stuck for a while. Help me step back and see it fresh.'' \\
\bottomrule
\end{tabular}
\end{table}

\section{Code Review Prompts}

Use when reviewing code or implementations.

\begin{table}[H]
\centering
\begin{tabular}{p{3cm}p{9cm}}
\toprule
\textbf{Situation} & \textbf{Prompt} \\
\midrule
Understanding & ``Explain this code as if I'm new to the team. What's the key insight?'' \\
Hidden assumptions & ``Where are the assumptions buried in this code?'' \\
Edge cases & ``What edge cases does this not handle?'' \\
Silent failures & ``Where could this fail silently?'' \\
Simplification & ``Is there a simpler way to achieve the same result?'' \\
\bottomrule
\end{tabular}
\end{table}

\section{Synthesis Prompts}

Use when wrapping up work or capturing learning.

\begin{table}[H]
\centering
\begin{tabular}{p{3cm}p{9cm}}
\toprule
\textbf{Situation} & \textbf{Prompt} \\
\midrule
Session summary & ``What's the one thing I should remember from this session?'' \\
Learning capture & ``What did I learn that I didn't know before?'' \\
Next steps & ``What question should I be asking next?'' \\
Pre-compaction & ``What context is essential for continuing this work?'' \\
\bottomrule
\end{tabular}
\end{table}

\section{Meta Prompts}

Use for reflection on the collaboration itself.

\begin{table}[H]
\centering
\begin{tabular}{p{3cm}p{9cm}}
\toprule
\textbf{Situation} & \textbf{Prompt} \\
\midrule
Engagement check & ``Am I extracting answers or building understanding?'' \\
Collaboration quality & ``What would make this collaboration more effective?'' \\
Calibration & ``I keep being surprised by X. How should I update my model?'' \\
\bottomrule
\end{tabular}
\end{table}

\section{Protocols and Templates}\label{sec:protocols-templates}\index{Protocols}

These structured protocols provide templates for specific DCF modes.

\subsection{Task State Markers (for /dcf decompose)}\index{Task State Markers}

Track task progress explicitly to prevent goal drift:

\begin{table}[H]
\centering
\begin{tabular}{ll}
\toprule
\textbf{State} & \textbf{Meaning} \\
\midrule
\texttt{brainstormed} & Identified but not validated \\
\texttt{validated} & Requirements confirmed with user \\
\texttt{blocked} & Waiting on dependency or decision \\
\texttt{in\_progress} & Currently being worked \\
\texttt{review} & Complete, awaiting verification \\
\texttt{done} & Verified and accepted \\
\bottomrule
\end{tabular}
\caption{Task State Markers for Decomposition}
\end{table}

Use these states in decomposition output to make progress visible. State transitions should be explicit: \texttt{brainstormed} $\rightarrow$ \texttt{validated} requires user confirmation; \texttt{in\_progress} $\rightarrow$ \texttt{review} requires implementation completion.

\subsection{Chain-of-Verification Protocol (for /dcf verify)}\index{Chain-of-Verification}

For high-stakes claims requiring rigorous verification:

\begin{enumerate}
    \item \textbf{State the claim}: What exactly are we verifying?
    \item \textbf{Identify verification questions}: 2-3 specific, answerable questions
    \item \textbf{Gather evidence independently}: For each question, find evidence, source, and confidence level
    \item \textbf{Synthesize}: Which questions agree? Conflict? What gaps remain?
    \item \textbf{Verdict}: Verified / Partially Verified / Unverified / Refuted, with confidence and caveats
\end{enumerate}

The key insight: gathering evidence independently for each question prevents confirmation bias. If you answer all questions from the same source, you're not verifying---you're echoing.

\subsection{Multi-Perspective Synthesis (for /dcf challenge)}\index{Multi-Perspective Synthesis}

For complex decisions, generate challenges from four frames before synthesis:

\begin{table}[H]
\centering
\begin{tabular}{lp{8cm}}
\toprule
\textbf{Frame} & \textbf{Question} \\
\midrule
Technical & Does this work mechanically? What could break? \\
User & Does this serve the actual need? What's the user experience? \\
Maintenance & Will future-you thank or curse present-you? \\
Adversarial & How would someone trying to break this succeed? \\
\bottomrule
\end{tabular}
\caption{Multi-Perspective Challenge Frames}
\end{table}

After generating each perspective's verdict (Support/Oppose/Neutral), synthesize:
\begin{itemize}
    \item \textbf{Converging concerns}: What multiple frames flag as problematic
    \item \textbf{Diverging assessments}: Where frames disagree (may indicate robustness)
    \item \textbf{Blind spots}: What no frame addressed
\end{itemize}

If concerns converge across frames, that's a critical vulnerability. If they diverge, the position may be stronger than initial analysis suggested.

